\documentclass[11pt,a4paper]{ctexart}
\usepackage[margin=2.2cm]{geometry}
\usepackage{xcolor}
\usepackage{enumitem}
\usepackage{amssymb}
\usepackage[normalem]{ulem}
\usepackage{titlesec}
\usepackage{hyperref}
\definecolor{navy}{HTML}{1A365D}
\definecolor{ink}{HTML}{1F2933}
\definecolor{muted}{HTML}{66788A}
\definecolor{soft}{HTML}{EEF4FA}
\definecolor{gold}{HTML}{8A5A12}
\hypersetup{colorlinks=true,linkcolor=navy,urlcolor=navy}
\setlist{nosep}
\titleformat{\section}{\Large\bfseries\color{navy}}{}{0pt}{}
\titleformat{\subsection}{\large\bfseries\color{navy}}{}{0pt}{}
\pagestyle{plain}

\title{\textbf{八下动词不定式专项}\\[4pt]
\large 人教版（新课纲版）八年级下册 · 动词不定式（Infinitives）}
\author{初中英语资料站}
\date{2026-11-12}
\begin{document}
\maketitle

\section*{使用说明}
\begin{itemize}[leftmargin=1.4em]
\item 本讲义对应人教版（新课纲版）八年级下册语法要求：动词不定式作宾语、目的状语、宾语补足语，以及否定式、疑问词结构和「不带 to」的情况。
\item 共 5 讲语法、6 级 48 题。
\item 在线版：\url{https://1901151458-prog.github.io/english-materials/g8b-infinitives/}
\end{itemize}

\section{构成与三种句法作用}
动词不定式的基本形式是「\textbf{to + 动词原形}」，没有人称和数的变化，不能作谓语，但可以有自己的宾语和状语。
\begin{center}
\begin{tabular}{p{3.0cm}p{11.0cm}}
\hline
\textbf{句法作用} & \textbf{例句}\\
\hline
作宾语 & Do you want \underline{to visit} the Great Wall? / I like \underline{to play} badminton in my free time.\\
作目的状语 & What do you do \underline{to pass} the time? / Allen ran to the living room \underline{to get} the fire extinguisher.\\
作宾语补足语 & My instructor encouraged me \underline{to get back up}. / Doctor Cooper asked him \underline{to be} more careful.\\
\hline
\end{tabular}
\end{center}

\section{否定式：not + 不定式}
动词不定式的否定形式由「\textbf{not + 动词不定式}」构成，not 放在 to 之前。
\begin{itemize}[leftmargin=1.4em]
\item I would prefer \underline{not to talk} about it.（× to not talk）
\item ``Try \underline{not to make} any noise when you eat,'' his mother said.
\item My parents told me \underline{not to push} myself too hard.
\end{itemize}

\section{疑问词 + 不定式}
动词不定式可以和 what、which、how、where、when 等疑问词连用，构成「疑问词 + 不定式」结构，相当于一个名词短语。
\begin{itemize}[leftmargin=1.4em]
\item He didn't know \underline{what to say} to Ella.
\item She asked me \underline{how to use} the washing machine.
\item The tour guide will tell you \underline{where to go} next.
\item I have no idea \underline{which book to choose}.
\end{itemize}

\section{不带 to 的不定式}
有一些动词后用作宾语补足语的不定式通常不\textbf{带 to}，主要有两类：
\begin{itemize}[leftmargin=1.4em]
\item \textbf{感官动词}：see、hear、watch、feel、notice\\
The referee saw Matt \underline{push} the other player. / He heard someone \underline{walk} down the stairway.
\item \textbf{使役动词}：let、make、have\\
Please let me \underline{know}. / Hobbies can make you \underline{feel} good about yourself.
\end{itemize}
\textbf{注意：}这两类动词\textbf{变被动语态时，to 要还原}：He was seen \underline{to push} the other player.

\section{作宾语 vs 作宾补}
\begin{itemize}[leftmargin=1.4em]
\item \textbf{作宾语：动词 + to do} —— want / decide / hope / agree / promise / plan / refuse / fail\\
She \underline{decided to study} harder.
\item \textbf{作宾补：动词 + sb + to do} —— tell / ask / want / encourage / advise / teach\\
She \underline{told me to study} harder.
\end{itemize}
\textbf{区别：}不定式前面有没有「人（宾语）」——有就是宾语补足语，没有就是宾语。


\newpage
\section{分级练习（48 题）}
\noindent 按难度分级排列，先做基础，再做成分与扩展，最后做综合题。

\subsection*{第 1 级 · 基础 · 不定式的形式}
\noindent\textcolor{gray}{to + 动词原形；没有时态和人称变化。}\par\smallskip
\begin{enumerate}[leftmargin=*,itemsep=4pt]
\item I want \_\_\_ a doctor when I grow up.
\begin{enumerate}[label=\Alph*.,leftmargin=2em,itemsep=1pt,topsep=2pt]
\item to be
\item being
\item be
\item been
\end{enumerate}
\item She decided \_\_\_ harder from now on.
\begin{enumerate}[label=\Alph*.,leftmargin=2em,itemsep=1pt,topsep=2pt]
\item study
\item to study
\item studying
\item studied
\end{enumerate}
\item They hope \_\_\_ the game tomorrow.
\begin{enumerate}[label=\Alph*.,leftmargin=2em,itemsep=1pt,topsep=2pt]
\item win
\item to win
\item winning
\item won
\end{enumerate}
\item He agreed \_\_\_ us with the project.
\begin{enumerate}[label=\Alph*.,leftmargin=2em,itemsep=1pt,topsep=2pt]
\item to help
\item helping
\item help
\item helped
\end{enumerate}
\item We plan \_\_\_ a trip next month.
\begin{enumerate}[label=\Alph*.,leftmargin=2em,itemsep=1pt,topsep=2pt]
\item take
\item to take
\item taking
\item took
\end{enumerate}
\item She refused \_\_\_ the question.
\begin{enumerate}[label=\Alph*.,leftmargin=2em,itemsep=1pt,topsep=2pt]
\item answer
\item to answer
\item answering
\item answered
\end{enumerate}
\item I’d like \_\_\_ something about the city.
\begin{enumerate}[label=\Alph*.,leftmargin=2em,itemsep=1pt,topsep=2pt]
\item know
\item to know
\item knowing
\item knew
\end{enumerate}
\item He failed \_\_\_ the exam last week.
\begin{enumerate}[label=\Alph*.,leftmargin=2em,itemsep=1pt,topsep=2pt]
\item pass
\item to pass
\item passing
\item passed
\end{enumerate}
\end{enumerate}

\subsection*{第 2 级 · 作宾语}
\noindent\textcolor{gray}{want / decide / hope / agree / promise 等动词后接不定式。}\par\smallskip
\begin{enumerate}[leftmargin=*,itemsep=4pt]
\item Do you want \_\_\_ the Great Wall?
\begin{enumerate}[label=\Alph*.,leftmargin=2em,itemsep=1pt,topsep=2pt]
\item to visit
\item visiting
\item visit
\item visits
\end{enumerate}
\item I like \_\_\_ badminton in my free time.
\begin{enumerate}[label=\Alph*.,leftmargin=2em,itemsep=1pt,topsep=2pt]
\item to play
\item play
\item played
\item plays
\end{enumerate}
\item On Saturday, James decided \_\_\_ fried chicken.
\begin{enumerate}[label=\Alph*.,leftmargin=2em,itemsep=1pt,topsep=2pt]
\item to make
\item making
\item make
\item made
\end{enumerate}
\item They hope \_\_\_ the problem by themselves.
\begin{enumerate}[label=\Alph*.,leftmargin=2em,itemsep=1pt,topsep=2pt]
\item solve
\item to solve
\item solving
\item solved
\end{enumerate}
\item I forgot \_\_\_ the door when I left.
\begin{enumerate}[label=\Alph*.,leftmargin=2em,itemsep=1pt,topsep=2pt]
\item lock
\item to lock
\item locking
\item locked
\end{enumerate}
\item She promised \_\_\_ on time.
\begin{enumerate}[label=\Alph*.,leftmargin=2em,itemsep=1pt,topsep=2pt]
\item arrive
\item to arrive
\item arriving
\item arrived
\end{enumerate}
\item We agreed \_\_\_ the meeting on Friday.
\begin{enumerate}[label=\Alph*.,leftmargin=2em,itemsep=1pt,topsep=2pt]
\item hold
\item to hold
\item holding
\item held
\end{enumerate}
\item He chose \_\_\_ the piano instead of the violin.
\begin{enumerate}[label=\Alph*.,leftmargin=2em,itemsep=1pt,topsep=2pt]
\item learn
\item to learn
\item learning
\item learned
\end{enumerate}
\end{enumerate}

\subsection*{第 3 级 · 作目的状语}
\noindent\textcolor{gray}{表示「为了做某事」，说明动作的目的。}\par\smallskip
\begin{enumerate}[leftmargin=*,itemsep=4pt]
\item What do you do \_\_\_ the time?
\begin{enumerate}[label=\Alph*.,leftmargin=2em,itemsep=1pt,topsep=2pt]
\item to pass
\item passing
\item pass
\item passed
\end{enumerate}
\item Allen ran to the living room \_\_\_ the fire extinguisher.
\begin{enumerate}[label=\Alph*.,leftmargin=2em,itemsep=1pt,topsep=2pt]
\item to get
\item getting
\item get
\item got
\end{enumerate}
\item Do not stand up \_\_\_ for food across the table.
\begin{enumerate}[label=\Alph*.,leftmargin=2em,itemsep=1pt,topsep=2pt]
\item to reach
\item reaching
\item reach
\item reached
\end{enumerate}
\item She went to the library \_\_\_ some books.
\begin{enumerate}[label=\Alph*.,leftmargin=2em,itemsep=1pt,topsep=2pt]
\item to borrow
\item borrowing
\item borrow
\item borrowed
\end{enumerate}
\item He got up early \_\_\_ the first bus.
\begin{enumerate}[label=\Alph*.,leftmargin=2em,itemsep=1pt,topsep=2pt]
\item to catch
\item catching
\item catch
\item caught
\end{enumerate}
\item We stayed up late \_\_\_ the report.
\begin{enumerate}[label=\Alph*.,leftmargin=2em,itemsep=1pt,topsep=2pt]
\item to finish
\item finishing
\item finish
\item finished
\end{enumerate}
\item She spoke loudly \_\_\_ everyone could hear her.
\begin{enumerate}[label=\Alph*.,leftmargin=2em,itemsep=1pt,topsep=2pt]
\item to make sure
\item making sure
\item make sure
\item made sure
\end{enumerate}
\item They went to the station \_\_\_ their friend.
\begin{enumerate}[label=\Alph*.,leftmargin=2em,itemsep=1pt,topsep=2pt]
\item to meet
\item meeting
\item meet
\item met
\end{enumerate}
\end{enumerate}

\subsection*{第 4 级 · 作宾语补足语}
\noindent\textcolor{gray}{tell / ask / want / encourage sb to do sth。}\par\smallskip
\begin{enumerate}[leftmargin=*,itemsep=4pt]
\item My instructor encouraged me \_\_\_ back up whenever I fell.
\begin{enumerate}[label=\Alph*.,leftmargin=2em,itemsep=1pt,topsep=2pt]
\item to get
\item getting
\item get
\item got
\end{enumerate}
\item My parents told me \_\_\_ myself too hard.
\begin{enumerate}[label=\Alph*.,leftmargin=2em,itemsep=1pt,topsep=2pt]
\item not to push
\item not push
\item to not push
\item pushing
\end{enumerate}
\item Doctor Cooper asked him \_\_\_ more careful in the future.
\begin{enumerate}[label=\Alph*.,leftmargin=2em,itemsep=1pt,topsep=2pt]
\item to be
\item be
\item being
\item been
\end{enumerate}
\item The teacher wants us \_\_\_ the text aloud.
\begin{enumerate}[label=\Alph*.,leftmargin=2em,itemsep=1pt,topsep=2pt]
\item to read
\item read
\item reading
\item readed
\end{enumerate}
\item She told the children \_\_\_ in the street.
\begin{enumerate}[label=\Alph*.,leftmargin=2em,itemsep=1pt,topsep=2pt]
\item not to play
\item not play
\item to not play
\item playing
\end{enumerate}
\item The doctor advised him \_\_\_ more exercise.
\begin{enumerate}[label=\Alph*.,leftmargin=2em,itemsep=1pt,topsep=2pt]
\item to take
\item take
\item taking
\item took
\end{enumerate}
\item Our teacher asked us \_\_\_ quiet in the library.
\begin{enumerate}[label=\Alph*.,leftmargin=2em,itemsep=1pt,topsep=2pt]
\item to keep
\item keep
\item keeping
\item kept
\end{enumerate}
\item The coach encouraged the players \_\_\_ their best.
\begin{enumerate}[label=\Alph*.,leftmargin=2em,itemsep=1pt,topsep=2pt]
\item to try
\item try
\item trying
\item tried
\end{enumerate}
\end{enumerate}

\subsection*{第 5 级 · 否定式与疑问词 + 不定式}
\noindent\textcolor{gray}{not to do；what/how/where/which + to do。}\par\smallskip
\begin{enumerate}[leftmargin=*,itemsep=4pt]
\item I would prefer \_\_\_ about it. It’s private.
\begin{enumerate}[label=\Alph*.,leftmargin=2em,itemsep=1pt,topsep=2pt]
\item not to talk
\item not talking
\item to not talk
\item talk not
\end{enumerate}
\item ’’Try \_\_\_ any noise when you eat,’’ his mother said.
\begin{enumerate}[label=\Alph*.,leftmargin=2em,itemsep=1pt,topsep=2pt]
\item not to make
\item not make
\item to not make
\item making
\end{enumerate}
\item He didn’t know \_\_\_ to Ella.
\begin{enumerate}[label=\Alph*.,leftmargin=2em,itemsep=1pt,topsep=2pt]
\item what to say
\item what say
\item how say
\item what saying
\end{enumerate}
\item She asked me \_\_\_ the washing machine.
\begin{enumerate}[label=\Alph*.,leftmargin=2em,itemsep=1pt,topsep=2pt]
\item how to use
\item how use
\item what use
\item how using
\end{enumerate}
\item The tour guide will tell you \_\_\_ next.
\begin{enumerate}[label=\Alph*.,leftmargin=2em,itemsep=1pt,topsep=2pt]
\item where to go
\item where go
\item where going
\item to where go
\end{enumerate}
\item I have no idea \_\_\_ for my book report.
\begin{enumerate}[label=\Alph*.,leftmargin=2em,itemsep=1pt,topsep=2pt]
\item which book to choose
\item which book choose
\item which to choose book
\item choose which book
\end{enumerate}
\item She was not sure \_\_\_ next.
\begin{enumerate}[label=\Alph*.,leftmargin=2em,itemsep=1pt,topsep=2pt]
\item what to do
\item what do
\item how do
\item what doing
\end{enumerate}
\item He asked me \_\_\_ him with his homework.
\begin{enumerate}[label=\Alph*.,leftmargin=2em,itemsep=1pt,topsep=2pt]
\item to help
\item help
\item helping
\item helped
\end{enumerate}
\end{enumerate}

\subsection*{第 6 级 · 不带 to 的不定式}
\noindent\textcolor{gray}{感官动词 see/hear/watch/feel/notice 与使役动词 let/make/have。}\par\smallskip
\begin{enumerate}[leftmargin=*,itemsep=4pt]
\item The referee saw Matt \_\_\_ the other player.
\begin{enumerate}[label=\Alph*.,leftmargin=2em,itemsep=1pt,topsep=2pt]
\item push
\item to push
\item pushing
\item pushed
\end{enumerate}
\item He heard someone \_\_\_ down the stairway.
\begin{enumerate}[label=\Alph*.,leftmargin=2em,itemsep=1pt,topsep=2pt]
\item walk
\item to walk
\item walked
\item walks
\end{enumerate}
\item If you have any questions, please let me \_\_\_.
\begin{enumerate}[label=\Alph*.,leftmargin=2em,itemsep=1pt,topsep=2pt]
\item know
\item to know
\item knowing
\item known
\end{enumerate}
\item Hobbies can make you \_\_\_ good about yourself.
\begin{enumerate}[label=\Alph*.,leftmargin=2em,itemsep=1pt,topsep=2pt]
\item feel
\item to feel
\item feeling
\item felt
\end{enumerate}
\item I noticed him \_\_\_ a quick look at his watch.
\begin{enumerate}[label=\Alph*.,leftmargin=2em,itemsep=1pt,topsep=2pt]
\item take
\item to take
\item taking
\item took
\end{enumerate}
\item She watched the children \_\_\_ in the garden.
\begin{enumerate}[label=\Alph*.,leftmargin=2em,itemsep=1pt,topsep=2pt]
\item play
\item to play
\item played
\item plays
\end{enumerate}
\item Don’t make the baby \_\_\_.
\begin{enumerate}[label=\Alph*.,leftmargin=2em,itemsep=1pt,topsep=2pt]
\item cry
\item to cry
\item crying
\item cried
\end{enumerate}
\item I felt someone \_\_\_ my shoulder.
\begin{enumerate}[label=\Alph*.,leftmargin=2em,itemsep=1pt,topsep=2pt]
\item touch
\item to touch
\item touched
\item touches
\end{enumerate}
\end{enumerate}

\newpage
\section*{答案与解析}
\small
\begin{enumerate}[leftmargin=1.6em,itemsep=2pt]
\item \textbf{I1-1 A}：want 后接不定式 to be。
\item \textbf{I1-2 B}：decide 后接不定式 to study。
\item \textbf{I1-3 B}：hope 后接不定式 to win。
\item \textbf{I1-4 A}：agree 后接不定式 to help。
\item \textbf{I1-5 B}：plan 后接不定式 to take。
\item \textbf{I1-6 B}：refuse 后接不定式 to answer。
\item \textbf{I1-7 B}：would like 后接不定式 to know。
\item \textbf{I1-8 B}：fail 后接不定式 to pass。
\item \textbf{I2-1 A}：want to visit；教材例句（表6）。
\item \textbf{I2-2 A}：like to do sth；教材例句（表6）。
\item \textbf{I2-3 A}：decide to make；教材例句（表6）。
\item \textbf{I2-4 B}：hope 后接不定式作宾语。
\item \textbf{I2-5 B}：forget to do 表示「忘记去做」（还没做）。
\item \textbf{I2-6 B}：promise 后接不定式 to arrive。
\item \textbf{I2-7 B}：agree 后接不定式 to hold。
\item \textbf{I2-8 B}：choose 后接不定式 to learn。
\item \textbf{I3-1 A}：to pass the time 表示目的「为了打发时间」；教材例句（表6）。
\item \textbf{I3-2 A}：不定式作目的状语；教材例句（表6）。
\item \textbf{I3-3 A}：不定式表目的；教材例句（表6）。
\item \textbf{I3-4 A}：去图书馆的目的是借书，用不定式表目的。
\item \textbf{I3-5 A}：早起的目的：赶上第一班车。
\item \textbf{I3-6 A}：熬夜的目的：完成报告。
\item \textbf{I3-7 A}：不定式表目的（这里也可用 so that 从句）。
\item \textbf{I3-8 A}：去车站的目的是接朋友。
\item \textbf{I4-1 A}：encourage sb to do；教材例句（表6）。
\item \textbf{I4-2 A}：tell sb not to do；否定式；教材例句（表6）。
\item \textbf{I4-3 A}：ask sb to do；教材例句（表6）。
\item \textbf{I4-4 A}：want sb to do sth。
\item \textbf{I4-5 A}：tell sb not to do sth 否定式。
\item \textbf{I4-6 A}：advise sb to do sth。
\item \textbf{I4-7 A}：ask sb to do sth。
\item \textbf{I4-8 A}：encourage sb to do sth。
\item \textbf{I5-1 A}：不定式的否定式：not to do；教材例句。
\item \textbf{I5-2 A}：try not to do sth；教材例句。
\item \textbf{I5-3 A}：疑问词 + 不定式：what to say；教材例句。
\item \textbf{I5-4 A}：how to use；教材例句。
\item \textbf{I5-5 A}：where to go；教材例句。
\item \textbf{I5-6 A}：which + 名词 + 不定式。
\item \textbf{I5-7 A}：what to do 表示「该做什么」。
\item \textbf{I5-8 A}：ask sb to do sth 作宾语补足语。
\item \textbf{I6-1 A}：感官动词 see sb do sth，不定式不带 to；教材例句。
\item \textbf{I6-2 A}：感官动词 hear sb do sth，不带 to；教材例句。
\item \textbf{I6-3 A}：使役动词 let sb do sth，不带 to；教材例句。
\item \textbf{I6-4 A}：使役动词 make sb do sth，不带 to；教材例句。
\item \textbf{I6-5 A}：感官动词 notice sb do sth，不带 to；教材例句。
\item \textbf{I6-6 A}：感官动词 watch sb do sth，不带 to。
\item \textbf{I6-7 A}：使役动词 make sb do sth，不带 to。
\item \textbf{I6-8 A}：感官动词 feel sb do sth，不带 to。
\end{enumerate}

\end{document}
